The natural branch uses the separate raw source of package 133, with its
dominant common digit modulus $M=Nq^{x+1}$. Its six exact powers and eleven modular numerator requests share fifteen
principal indices: four power indices, one exact-binomial upper index,
nine CRT representatives, and one constructor label. All raw coordinates in this
branch are nonnegative. The new linear beta modulus of
Section~\ref{sec:cofinal} is combined with a single bound-and-congruence
row; the one-row mechanism is credited to package 128.

Let $\mathcal S$ be the twelve distinct actual bases, exponents, and outputs
of the six power requests. Put
\[
 B_*=2+\sum\mathcal S+\sum_{i=1}^9a_i+J_{\rm label},
 \qquad W=2B_*.
\]
Every requested index $e_i$ is positive and less than $B_*$ directly from
the raw definitions. This assertion uses no power or factorial semantics.
The retained lower frame has
\[
 L^2-\bigl((W+1)^2-1\bigr)(WZ)^2=1,\qquad Z=z_0+1,
 \qquad A=L+\sigma+1,\quad \sigma\ge0.
\]
The shift by one is an existence-preserving restriction: the later common
root congruence admits cofinally many compatible $A$. The lower frame
gives $L\ge X_W(W+1)\ge(W+1)^W$ and $A-1\ge L>W^3$.

Retain the complete indexed anchor at $W$, with ordinate $C$, and replace
the old separate cap and the two staged residues by the single quadratic
row
\begin{equation}\label{eq:naturalcap}
 \boxed{C=W+\sum_{i=1}^{15}(Y_i-e_i)^2+(A-1)Q,\qquad Q\ge0.}
\end{equation}
The nonnegative domain is essential here. It first supplies $C\ge W$,
so the anchor theorem yields $C=Y_W(A)$ before any seed-index bound is
invoked.

\begin{lemma}[One-row identification of all indices]
Under these raw natural conditions, \eqref{eq:naturalcap} forces
$Y_i=Y_{e_i}(A)$ for every principal pair. Conversely, canonical pairs
admit a positive integer $Q$ satisfying that row.
\end{lemma}
\begin{proof}
Write $B=B_*$ temporarily. If $1\le e<B$, put $z=Y_B(A)$ and
$x=X_B(A)$. The norm gives $(A-1)z<x\le Az$. Since $z\ge B$,
\[
 Y_{B+1}(A)-e>(2A-2)z+1,\qquad Y_{2B}(A)=2xz\le2Az^2.
\]
For $A\ge2$, $(2A-2)^2\ge2A$, so
$(Y_{B+1}-e)^2>Y_{2B}=C$. But \eqref{eq:naturalcap} makes each squared
error smaller than $C$. Thus the actual index $j_i$ is at most $B$.
Reducing modulo $A-1$ now gives
\[
 A-1\mid\sum_i(j_i-e_i)^2,
 \qquad 0\le\sum_i(j_i-e_i)^2\le15(B-1)^2<W^3<A-1.
\]
Every error is zero.

For the converse choose all pairs canonically and set $C=Y_{2B}(A)$.
With $z=Y_{B-1}(A)$, the sum of the squared errors is at most $15z^2$,
whereas $C>2(A-1)A^2z^2$. Since $A>W+1$, their difference exceeds $W$.
Both $C-W$ and each $Y_{e_i}-e_i$ are divisible by $A-1$, so the
quotient in \eqref{eq:naturalcap} is a positive integer.
\end{proof}

All six powers are now exact by the retained small-residue identities;
the dominant common modulus, factorial fixed point, range digits, and
calculus decoding follow in the original order. In completeness, the
semantic data and all frame atoms are fixed first. The common-root
congruence then supplies a sufficiently large compatible $A$, followed
by the anchor and canonical seeds. There is no self-referential choice
of the frame.

For the linear beta restriction, use a new nonnegative coordinate
$v=V-1$ and the row $y=v$, replacing the old endpoint slack.
Soundness restores that slack as $y-F-1\ge0$ after the early ceiling
bootstrap; completeness is Lemma~\ref{lem:linear}. All occurrences of
$V$ and $A$ are shifted together in the literal graph.

Monic projection with explicit nonnegativity certificates gives the
235-variable equation with 184 residuals, 22,185 monomials, and height
90,911,633. A sparser 236-variable companion keeps two staged rows:
the first authenticates the five raw-bounded indices; only after exact
power and factorial recovery does the second authenticate the remaining
ten. It has 185 residuals, 15,645 monomials, and height 88,233,865.
Both have degree four and the same ZFC-provable consistency equivalence
over $\N$. Their complete literal residuals, expansions, inverse
projection certificates, and ordinary proofs are supplied in
\filepath{natural235/}.
